% Môn: Vật lý
% Lớp: 11
% Chương: Chương I. Dao động
% Bài: L11C1 Bài 2. Mô tả dao động điều hòa · Xác định pha ban đầu, pha dao động và trạng thái li độ tại thời điểm t
% ===== Câu 48 =====
% ID: L11C1B2-193-DS
% Mức: TH
\dangbt{Xác định pha ban đầu, pha dao động và trạng thái li độ tại thời điểm t}
\begin{ex}
% ID: L11C1B2-193-DS
% Mức: TH
Vật dao động theo $x=4\cos(2\pi t/1+\pi/3)$ cm. Xác định pha, li độ và chiều chuyển động tại $t=0$.
\choiceTF

\end{ex}

% ===== Câu 62 =====
% ID: L11C1B2-194-TLN
% Mức: TH
\begin{ex}
% ID: L11C1B2-194-TLN
% Mức: TH
Hai dao động cùng tần số có pha ban đầu $\pi/6$ và $\pi/2.$ Độ lệch pha bằng bao nhiêu lần $\pi$?
\shortans{01/03/2026}
\end{ex}

% ===== Câu 66 =====
% ID: L11C1B2-195-DS
% Mức: TH
\begin{ex}
% ID: L11C1B2-195-DS
% Mức: TH
Vật dao động theo $x=4\cos(2\pi t/1+\pi/3)$ cm. Xác định tính đúng sai.
\choiceTF
{\True Pha ban đầu là $\pi/3$ rad.}
{\True Tại $t=0$, li độ bằng $2\,\text{cm}$.}
{\True Chu kì là $1\,\text{s}$.}
{Tại $t=0$, vật chuyển động theo chiều dương.}
\end{ex}

% ===== Câu 111 =====
% ID: L11C1B2-196-TN
% Mức: NB
\begin{ex}
% ID: L11C1B2-196-TN
% Mức: NB
Phương trình dao động của một chất điểm có dạng $x = 4\cos\left(2\pi t - \dfrac{\pi}{3}\right)$ (cm). Pha ban đầu của dao động này là:
\choice
{\True $-\dfrac{\pi}{3}$ rad}
{$\dfrac{\pi}{3}$ rad}
{$2\pi$ rad}
{$-\dfrac{2\pi}{3}$ rad}
\loigiai{
• Đồng nhất phương trình dao động của vật $x = 4\cos\left(2\pi t - \dfrac{\pi}{3}\right)$ với phương trình tổng quát có dạng $x = A\cos(\omega t + \varphi)$.
 
• Nhận thấy phương trình đã ở dạng chuẩn dương với biên độ $A = 4$ cm $\gt  0$ và tần số góc $\omega = 2\pi$ rad/s $\gt  0$.
 
• Đối chiếu hệ số tự do trong pha dao động ta thu được pha ban đầu của dao động là $\varphi = -\dfrac{\pi}{3}$ rad.
}
\end{ex}

% ===== Câu 112 =====
% ID: L11C1B2-197-TN
% Mức: TH
\begin{ex}
% ID: L11C1B2-197-TN
% Mức: TH
Một chất điểm dao động điều hòa theo phương trình có dạng chưa chuẩn $x = -6\cos\left(5\pi t + \dfrac{\pi}{6}\right)$ (cm). Pha ban đầu của dao động này sau khi đã được chuẩn hóa về dạng phương trình có biên độ dương là:
\choice
{$\dfrac{\pi}{6}$ rad}
{\True $-\dfrac{5\pi}{6}$ rad}
{$-\dfrac{\pi}{6}$ rad}
{$\dfrac{5\pi}{6}$ rad}
\loigiai{
• Sử dụng hệ thức lượng giác để khử dấu trừ trước hàm cosin nhằm đưa phương trình về dạng chuẩn có biên độ dương: $-\cos\alpha = \cos(\alpha - \pi)$.
 
• Ta thực hiện biến đổi phương trình dao động: $x = -6\cos\left(5\pi t + \dfrac{\pi}{6}\right) = 6\cos\left(5\pi t + \dfrac{\pi}{6} - \pi\right) = 6\cos\left(5\pi t - \dfrac{5\pi}{6}\right)$ (cm).
 
• Đồng nhất phương trình thu được với phương trình tổng quát $x = A\cos(\omega t + \varphi)$ để suy ra pha ban đầu: $\varphi = -\dfrac{5\pi}{6}$ rad.
}
\end{ex}

% ===== Câu 113 =====
% ID: L11C1B2-198-TN
% Mức: VD
\begin{ex}
% ID: L11C1B2-198-TN
% Mức: VD
Một chất điểm dao động điều hòa với phương trình $x = 5\cos\left(10\pi t + \dfrac{\pi}{3}\right)$ (cm). Li độ của chất điểm tại thời điểm pha dao động bằng $\pi$ là:
\choice
{$5$ cm}
{\True $-5$ cm}
{$2,5$ cm}
{$-2,5$ cm}
\loigiai{
• Pha của dao động tại thời điểm $t$ được xác định bằng biểu thức ở trong hàm cosin: $\Phi = 10\pi t + \dfrac{\pi}{3}$.
 
• Đề bài cho biết tại thời điểm khảo sát pha dao động của vật đạt giá trị $\Phi = \pi$ rad.
 
• Thay trực tiếp giá trị của pha dao động vào phương trình li độ để tính li độ của chất điểm: $x = 5\cos(\pi) = 5 \cdot (-1) = -5$ cm.
}
\end{ex}

% ===== Câu 114 =====
% ID: L11C1B2-199-TN
% Mức: VD
\begin{ex}
% ID: L11C1B2-199-TN
% Mức: VD
Một chất điểm dao động điều hòa có phương trình $x = 5\sqrt{3}\cos\left(10\pi t + \dfrac{\pi}{3}\right)$ (cm). Tại thời điểm $t = 1$ s, tọa độ của chất điểm đạt giá trị bằng:
\choice
{$2,5$ cm}
{$-5\sqrt{3}$ cm}
{\True $2,5\sqrt{3}$ cm}
{$-2,5\sqrt{3}$ cm}
\loigiai{
• Thay giá trị thời điểm $t = 1$ s vào biểu thức li độ: $x = 5\sqrt{3}\cos\left(10\pi \cdot 1 + \dfrac{\pi}{3}\right) = 5\sqrt{3}\cos\left(10\pi + \dfrac{\pi}{3}\right)$ (cm).
 
• Do tính chất tuần hoàn tuần hoàn của hàm số cosin với chu kì $2\pi$, ta có: $\cos\left(10\pi + \dfrac{\pi}{3}\right) = \cos\left(\dfrac{\pi}{3}\right)$.
 
• Thực hiện tính toán giá trị: $x = 5\sqrt{3}\cos\left(\dfrac{\pi}{3}\right) = 5\sqrt{3} \cdot 0,5 = 2,5\sqrt{3}$ cm.
}
\end{ex}

% ===== Câu 115 =====
% ID: L11C1B2-200-TN
% Mức: VD
\begin{ex}
% ID: L11C1B2-200-TN
% Mức: VD
Một chất điểm dao động điều hòa với phương trình $x = 10\cos\left(\pi t + \dfrac{\pi}{2}\right)$ (cm). Li độ của chất điểm tại thời điểm $t = 1$ s là:
\choice
{\True $0$ cm}
{$10$ cm}
{$-10$ cm}
{$5$ cm}
\loigiai{
• Thay trực tiếp giá trị thời gian $t = 1$ s vào biểu thức của phương trình dao động điều hòa.
 
• Ta thu được phương trình: $x = 10\cos\left(\pi \cdot 1 + \dfrac{\pi}{2}\right) = 10\cos\left(\dfrac{3\pi}{2}\right)$ (cm).
 
• Tính toán giá trị lượng giác bằng máy tính cầm tay ở chế độ radian: $\cos\left(\dfrac{3\pi}{2}\right) = 0 \Rightarrow x = 0$ cm.
}
\end{ex}
